\documentclass{article}
\usepackage{/home/user1/code/tex/sty}
\usepackage{amsfonts}
\usepackage{amsmath}
\usepackage{geometry}
\usepackage{hyperref}
\setlength{\parindent}{0pt}
\geometry{top=1in,bottom=1in,left=1in,right=1in}
\begin{document}
\begin{center}
\textbf{\Large{Problem Set 1}}\\~\\
\begin{tabular}{l l}
Student&Agustin Romero\\
Email&ar1@stanford.edu\\
Instructor&D. Stanford\\
Assistants&G. Batra\\
Course&Physics 331\\
Institution&Stanford University\\
Date&\today\\
\end{tabular}
\end{center}
\section*{Problem 1.}
C.f. Peskin Ch. 9, Srednicky Ch. 44. We will calculate both (3.1) and (3.2) just to be thorough. For the c-number Gaussian integral,
\n{intf}{\int_{-\infty}^\infty}
\n{nih}{-\br{1}{2}}
\n{egh}{\br{1}{2}}
\n{dui}{(\det U)^{-1}}
\n{pin}{\Pi_{\alpha=1}^{\br{1}{2}n}}
\g{\intf \intf e^{\nih x^T A x} dx_1 ... dx_n &= \int e^{\nih y^T D y} dy_1 ... dy_n \\
&= \int dy_1...dy_n e^{\nih \sum lambda_i y_i^2}\\
\intertext{The $n$ integrals can be separated.}
&= \Pi_i \int dy_i e^{\nih \lambda_i y_i^2} \\
\intertext{The integrals are each Gaussian integrals.}
&= \Pi_i (\br{2\pi}{\lambda_i})^{1/2}
\intertext{Using $\Pi_i \lambda_i = \det D = \det(C^{-1} A C) = \det A$ we have the answer for the real valued case.}
&= (\br{(2\pi)^n}{\det A})^{1/2}}
Now for the Grassman-Gauss integral. Let $M$ be a complex antisymmetric matrix.
\g{\psi_i = U_{ij} \psi_j'}
\g{\int d^d\psi e^{\egh \psi^T M \psi} &= \dui \pin d^2 \psi_\alpha e^{\egh \psi^T M_\alpha \psi}\intertext{Use $\int d^2 \psi_\alpha e^{\egh \psi^T M_\alpha \psi = m_\alpha}$}
&= \dui \pin m_\alpha\\
&= \boxed{(\det M)^{\br{1}{2}}}}
\section*{Problem 2.}
C.f. Tong's SUSY notes.
\n{pneu}{\prod_{n=1}^\infty}
\n{sneu}{\sum_{n=1}^\infty}
\g{\zeich{nz8cb}\prod_{n=\br{1}{2},\br{3}{2},...}^\infty \br{n}{\beta} &= \pneu \br{2n}{\beta} = (\pneu \br{2}{\beta})(\pneu n)}
These two infinite products can be individually computed using zeta function regularization. For the first, let $\br{2}{\beta}=b$,
\g{\pneu \br{2}{\beta} &= \pneu b = e^{\sum_{n=1}^\infty \log b} = e^{\log b + \log b + ...} \\
&= e^{\log \beta (1 + 1 + ...)} 
\intertext{Use the regularized sum $\sum_{i=1}^\infty 1 = -\br{1}{2}$}
&= e^{-\br{1}{2}\log \beta}\\
&= \br{1}{\sqrt{\beta}}}
If $\beta>0$, then the answer is real, otherwise imaginary. Similarly for the other infinite product of $n$,
\g{\zeich{mcia}\pneu n = 1 \cdot 2 \cdot ... = e^{\sneu \log n} = e^{\log 1 + \log 2 + ...}}
To compute \gzeich{mcia} we will need $\sneu \log n$ which we will do quickly here,
\g{\zeta(s)=\sneu \br{1}{n^s}}
\g{\dd{x} a^x = a^x \log a}
\g{\dd{s}(n^{-s})=-n^{-s}\log n}
\g{\zeta(s)'=-\sneu \br{\log n}{n^s}}
\g{\zeta(0)'=-\sneu \log{n} = -\br{1}{2}\log(2\pi)}
So the infinite product of $n$ \gzeich{mcia} is
\g{\pneu n = e^{\sneu \log n}=e^{\br{1}{2}\log(2\pi)}=(2\pi)^{1/2}}
The answer \gzeich{nz8cb} is then,
\g{\prod_{\br{1}{2},\br{3}{2},...}^\infty \br{n}{\beta} &=(\br{2}{\beta})^{-\br{1}{2}} (2\pi)^{\br{1}{2}} = \boxed{\sqrt{\pi\beta}}}
\end{document}
